\chapter{Markov Chains and Queues}\label{ch:qm:markov-chains-and-queues}

At the best bid of a large-tick future sit 800 lots; at the best ask, 200. A market maker has just
joined the back of the bid queue and wants to know two things: how likely its order is to be filled
before the ask queue is eaten and the price ticks up, leaving the order stranded; and how long it
will wait. Both are questions about the time a random walk takes to hit zero. If orders join and
leave a queue at random times, each queue is a \omterm{def:qm:markov-chains-and-queues:ctmc}{Markov chain}, and the answers follow from the
chain's generator: 73\% and 82 seconds, for the rates of this chapter, against 99\% if the ask queue
were twice as long. This chapter sets out \omterm{def:qm:markov-chains-and-queues:ctmc}{continuous-time Markov chains}, the birth--death processes
and simple queues built from them, the \omterm{met:qm:markov-chains-and-queues:firststep}{first-step analysis} that computes hitting probabilities and
times, and a model of the two best queues of an order book in the spirit of Cont and de Larrard
(2013).

\section{Continuous-time Markov chains}

\begin{definition}[Markov chain, continuous-time Markov chain]\label{def:qm:markov-chains-and-queues:ctmc}
A \emph{Markov chain}\index{Markov chain} is a \omterm{def:qm:brownian-motion:markov}{Markov process} with a countable state space, observed
at integer times, with transition matrix $P_{ij} = \P(X_{n+1} = j \mid X_n = i)$. A
\emph{continuous-time Markov chain}\index{continuous-time Markov chain} is a right-continuous \omterm{def:qm:brownian-motion:markov}{Markov
process} $(X_t)_{t\ge0}$ with a countable state space whose transition probabilities $P_{ij}(t) = \P(X_{s+t}
= j \mid X_s = i)$ do not depend on $s$.
\end{definition}

\begin{definition}[Generator matrix, jump chain]\label{def:qm:markov-chains-and-queues:generator}
The \emph{generator matrix}\index{generator matrix} of a continuous-time chain is $\mathcal L = (q_{ij})$ with $q_{ij}
= \lim_{t\downarrow0}P_{ij}(t)/t \ge 0$ for $j \ne i$, the rate of jumps from $i$ to $j$, and $q_{ii} = -\sum_{j\ne i}q_{ij}$. The
\emph{jump chain}\index{jump chain} is the discrete-time chain of the states visited, with transition
probabilities $q_{ij}/q_i$, where $q_i = -q_{ii}$.
\end{definition}

\begin{proposition}[Construction and Kolmogorov equations]\label{prop:qm:markov-chains-and-queues:kolmogorov}
A finite chain with generator $\mathcal L$ stays in state $i$ for an exponential time with rate $q_i$, then moves
according to the \omterm{def:qm:markov-chains-and-queues:generator}{jump chain}, independently of the past. Its transition matrix is $P(t) = e^{t\mathcal L}$, which
solves $P'(t) = P(t)\mathcal L = \mathcal LP(t)$, the forward and backward equations.
\end{proposition}

\begin{proof}
The exponential holding time is the only memoryless law, which the Markov property imposes on the time
spent in a state. For small $h$, $P(t + h) = P(t)P(h) = P(t)(I + h\mathcal L + o(h))$; the matrix exponential is the
unique solution with $P(0) = I$.
\end{proof}

A continuous-time chain is the discrete generator of chapter~4 for a finite state space: the same
backward and forward equations, with a matrix instead of a differential operator. Its \omterm{def:qm:stochastic-differential-equations:kolmogorov}{stationary
distribution} solves $\pi\mathcal L = 0$.

\begin{example}[The spread in ticks]\label{ex:qm:markov-chains-and-queues:spread}
Model a large-tick contract's spread as 1, 2 or 3 ticks, widening from 1 to 2 at rate 2 per second and from 2
to 3 at rate 1, narrowing from 2 to 1 at rate 5, from 3 to 2 at rate 4 and from 3 to 1 at rate 0.5. Solving $\pi\mathcal L
= 0$ gives $\pi = (0.676, 0.265, 0.059)$, a mean spread of 1.38 ticks; the spread sits at one tick for 0.5 seconds on
average before widening.
\end{example}

\begin{definition}[Detailed balance]\label{def:qm:markov-chains-and-queues:balance}
A distribution $\pi$ satisfies \emph{detailed balance}\index{detailed balance} with a generator if $\pi_iq_{ij} =
\pi_jq_{ji}$ for all $i, j$: in stationarity every transition is as frequent as its reverse, and the chain run
backwards in time has the same law (it is reversible).
\end{definition}

\omterm{def:qm:markov-chains-and-queues:balance}{Detailed balance} implies $\pi\mathcal L = 0$ by summing over $i$; the converse fails when probability circulates
around a cycle, as it does in the spread example (\cref{exo:qm:markov-chains-and-queues:4}). Birth--death
chains, which move one step at a time, always satisfy it.

\section{Birth--death processes and simple queues}

\begin{definition}[Birth--death process]\label{def:qm:markov-chains-and-queues:bd}
A \emph{birth--death process}\index{birth--death process} is a continuous-time chain on $\{0, 1, 2, \dots\}$ whose only
jumps are $n \to n + 1$ at rate $\lambda_n$ and $n \to n - 1$ at rate $\mu_n$.
\end{definition}

\begin{proposition}[Stationary law of a birth--death process]\label{prop:qm:markov-chains-and-queues:bdstat}
If $\sum_n\prod_{k=1}^n\lambda_{k-1}/\mu_k < \infty$, the stationary law is $\pi_n \propto \prod_{k=1}^n\lambda_{k-1}/\mu_k$.
\end{proposition}

\begin{proof}
\omterm{def:qm:markov-chains-and-queues:balance}{Detailed balance} between $n - 1$ and $n$ reads $\pi_{n-1}\lambda_{n-1} = \pi_n\mu_n$; iterate and normalise.
\end{proof}

\begin{definition}[M/M/1 queue]\label{def:qm:markov-chains-and-queues:mm1}
The \emph{M/M/1 queue}\index{M/M/1 queue} has Poisson arrivals at rate $\lambda$, exponential service times with
rate $\mu$ and one server; the number in the system is a \omterm{def:qm:markov-chains-and-queues:bd}{birth--death process} with $\lambda_n = \lambda$ and $\mu_n =
\mu$. (Kendall's notation: Markov arrivals, Markov service, one server.)
\end{definition}

With utilisation $\rho = \lambda/\mu < 1$ the stationary law is geometric, $\pi_n = (1 - \rho)\rho^n$, so the mean number in the
system is $L = \rho/(1 - \rho)$. Little's law, $L = \lambda W$ for the mean time $W$ spent in the system, holds for
almost any stationary queue (Little, 1961), and gives $W = 1/(\mu - \lambda)$. Both explode as $\rho \to 1$: a matching
engine or a risk gate that runs at 90\% of its capacity has nine messages waiting on average, at 95\%
nineteen (\cref{fig:qm:markov-chains-and-queues:mm1}). One Quant Book~13 sizes its queues with these
formulas.

\begin{omfigure}
\begin{tikzpicture}[font=\small]
\foreach \i/\x in {0/0,1/2.2,2/4.4,3/6.6} {\node[draw=omDef,circle,thick,minimum size=0.85cm,fill=omDef!6] (s\i) at (\x,0) {$\i$};}
\node (dots) at (8.6,0) {$\cdots$};
\foreach \a/\b/\l in {0/1/0,1/2/1,2/3/2} {
  \draw[-{Stealth},thick,omThm] (s\a) to[bend left=35] node[above] {$\lambda_{\l}$} (s\b);
  \pgfmathtruncatemacro{\m}{\l+1}
  \draw[-{Stealth},thick,omDef] (s\b) to[bend left=35] node[below] {$\mu_{\m}$} (s\a);}
\draw[-{Stealth},thick,omThm] (s3) to[bend left=35] node[above] {$\lambda_3$} (dots);
\draw[-{Stealth},thick,omDef] (dots) to[bend left=35] node[below] {$\mu_4$} (s3);
\end{tikzpicture}

\omcaption{A birth--death chain: from state $n$ the only moves are up at rate $\lambda_n$ and down at rate $\mu_n$. A best
queue counted in orders is one, with limit orders as births and market orders and cancellations as
deaths.}\label{fig:qm:markov-chains-and-queues:bd}
\end{omfigure}

\begin{omfigure}
\begin{tikzpicture}
\begin{axis}[width=11cm,height=4.8cm,
  xlabel={utilisation $\rho = \lambda/\mu$},ylabel={mean number in system},
  tick label style={font=\small},label style={font=\small},
  xmin=0,xmax=0.95,ymin=0,ymax=10.5,grid=major,grid style={gray!25},
  legend columns=3,legend style={font=\footnotesize,at={(0.5,-0.3)},anchor=north,draw=none,fill=none,/tikz/every even column/.append style={column sep=8pt}},
  table/col sep=comma]
\addplot[omDef,very thick] table[x=rho,y=theory]{figdata/methods/08-markov-chains-and-queues/mm1.csv};
\addplot[omThm,only marks,mark=*,mark size=1.8pt] table[x=rho,y=sim]{figdata/methods/08-markov-chains-and-queues/mm1.csv};
\addplot[black,only marks,mark=x,mark size=2.6pt] table[x=rho,y=lamW]{figdata/methods/08-markov-chains-and-queues/mm1.csv};
\legend{$\rho/(1-\rho)$,time average (sim.),$\lambda W$ (sim.)}
\end{axis}
\end{tikzpicture}

\omcaption{The \omterm{def:qm:markov-chains-and-queues:mm1}{M/M/1 queue}: mean number in the system against utilisation, from the stationary law and from
200\,000 simulated customers per point, whose time-averaged count agrees with $\lambda W$ (Little's law). Near
saturation the simulation converges slowly. Data: the chapter's tutorial,
seeded.}\label{fig:qm:markov-chains-and-queues:mm1}
\end{omfigure}

\section{Hitting probabilities and hitting times}

\begin{definition}[Absorbing state, hitting probability]\label{def:qm:markov-chains-and-queues:hitting}
A state is an \emph{absorbing state}\index{absorbing state} if the chain never leaves it ($q_i = 0$). For a set $A$ of
states, the \emph{hitting probability}\index{hitting probability} $h_i = \P_i(\text{the chain ever enters } A)$.
\end{definition}

\begin{method}[First-step analysis]\label{met:qm:markov-chains-and-queues:firststep}
\emph{First-step analysis}\index{first-step analysis} conditions on the first jump. Hitting probabilities solve
$h_i = 1$ on $A$ and $\sum_jq_{ij}h_j = 0$ off $A$; expected \omterm{def:qm:brownian-motion:passage}{hitting times} $m_i$ solve $m_i = 0$ on $A$ and $\sum_jq_{ij}m_j =
-1$ off $A$ (the minimal nonnegative solutions). For a birth--death chain these systems are tridiagonal and
cost $O(n)$ to solve.
\end{method}

For a queue with constant rates $\lambda$ and $\mu$, the probability of reaching $n$ before 0 from $i$ is the
gambler's ruin of chapter~1, $(1 - (\mu/\lambda)^i)/(1 - (\mu/\lambda)^n)$, and the expected time to empty from $i$ when $\mu > \lambda$
is $i/(\mu - \lambda)$. Means hide a lot here: the law of the emptying time is skewed, and a race between two
queues needs the whole law. It comes from uniformisation: with $L$ at least every total rate, the chain
is a \omterm{def:qm:markov-chains-and-queues:generator}{jump chain} with transition matrix $I + \mathcal L/L$ run at the times of a \omterm{def:qm:jump-processes:poisson}{Poisson process} of rate $L$, so
$P(t) = \sum_k e^{-Lt}\frac{(Lt)^k}{k!}(I + \mathcal L/L)^k$, a sum of positive terms that is stable to compute
(tutorial, first listing).

\section{Queues sized for order books}

Count each best queue in orders of 10 lots. On each side, limit orders join at rate $\lambda = 0.9$ per second, and
orders leave, by market orders ($0.6$) and cancellations ($0.4$), at rate $\nu = 1$: the queues drift toward
zero, and when one empties the price moves. Take the two queues independent. Then the mid price moves up
next if the ask queue's emptying time $T_a$ comes before the bid's, and
\[
\P(\text{up}) = \P(T_a < T_b) = \int_0^\infty\bigl(1 - F_b(t)\bigr)\,dF_a(t),
\]
with $F_a, F_b$ the laws of the emptying times. \Cref{fig:qm:markov-chains-and-queues:moveup} shows it against
the queue sizes: with 800 lots bid and 200 offered, the next move is up with probability 95.4\%. Queue
imbalance predicts the next tick, a fact One Quant Book~10 turns into a signal.

For the market maker's order at the back of the bid, with $k$ orders ahead, the fill time is the time for
the $k$ orders ahead to leave (at rate $\nu$ each) plus the wait for a market order (rate $0.6$), independent of
the ask queue. The fill probability is the race of that time against $T_a$
(\cref{fig:qm:markov-chains-and-queues:fill}). With 80 orders ahead and 20 on the ask, it is 73.2\%, and
73.4\% in 20\,000 simulated races; with 20 ahead it would be 99.3\%. A place near the front of the queue is
worth money, which is why time priority is contested so hard (One Quant Book~1, chapter~19).

\begin{omfigure}
\begin{tikzpicture}
\begin{groupplot}[group style={group size=2 by 1,horizontal sep=1.6cm},
  width=7.6cm,height=4.9cm,tick label style={font=\small},label style={font=\small},
  grid=major,grid style={gray!25},ymin=0,ymax=1.05,table/col sep=comma]
\nextgroupplot[xlabel={bid queue (orders of 10 lots)},ylabel={P(next move up)},xmin=0,xmax=100,
  legend style={font=\footnotesize,at={(0.97,0.03)},anchor=south east,fill=white,draw=none}]
\addplot[omDef,very thick] table[x=bid,y=ask10]{figdata/methods/08-markov-chains-and-queues/moveup.csv};
\addplot[omThm,very thick] table[x=bid,y=ask20]{figdata/methods/08-markov-chains-and-queues/moveup.csv};
\addplot[omProp,very thick] table[x=bid,y=ask40]{figdata/methods/08-markov-chains-and-queues/moveup.csv};
\legend{ask 10,ask 20,ask 40}
\nextgroupplot[xlabel={orders ahead of ours},ylabel={P(fill before the ask empties)},ylabel style={font=\footnotesize},xmin=0,xmax=120,
  legend style={font=\footnotesize,at={(0.03,0.03)},anchor=south west,fill=white,draw=none}]
\addplot[omDef,very thick] table[x=ahead,y=ask10]{figdata/methods/08-markov-chains-and-queues/fill.csv};
\addplot[omThm,very thick] table[x=ahead,y=ask20]{figdata/methods/08-markov-chains-and-queues/fill.csv};
\addplot[omProp,very thick] table[x=ahead,y=ask40]{figdata/methods/08-markov-chains-and-queues/fill.csv};
\addplot[black,only marks,mark=o,mark size=2.6pt] coordinates {(80,0.7324)};
\legend{ask 10,ask 20,ask 40}
\end{groupplot}
\end{tikzpicture}

\omcaption{Two best queues as independent birth--death chains (joins at 0.9, departures at 1 per second).
Left: probability that the mid price's next move is up, against the bid queue, for three ask queues. Right:
probability that an order with a given number of orders ahead of it at the bid is filled before the ask
queue empties; the circle marks the weekend problem, 80 ahead and 20 offered: 73\%. Data: the chapter's
tutorial.}\label{fig:qm:markov-chains-and-queues:moveup}\label{fig:qm:markov-chains-and-queues:fill}
\end{omfigure}

\section{Tutorial: the back of the queue}\label{tut:qm:markov-chains-and-queues:queue}

\begin{tutorial}
\textbf{Goal.} Compute the law of a queue's emptying time by uniformisation, race it against a fill time, and
check the result by simulation. \textbf{End state:}
\cref{fig:qm:markov-chains-and-queues:mm1,fig:qm:markov-chains-and-queues:moveup} and the probability 73\%.
\begin{tutsteps}
\item \textbf{The emptying time} of a queue, and the race between two independent times.
\omcode{code/firm/queues/firm_queues.py}{104}{128}{Uniformisation of a birth--death queue absorbed at zero, and the race of two times.}\label{lst:qm:markov-chains-and-queues:depletion}
\item \textbf{The fill time}: $k$ departures ahead, then a market order.
\omcode{code/firm/queues/firm_queues.py}{131}{149}{The fill time of an order at the back of the queue.}\label{lst:qm:markov-chains-and-queues:fill}
\item \textbf{Run} \texttt{problem()}, \texttt{spread\_chain()}, \texttt{mm1\_simulation(rho)} and
\texttt{fig\_queues.py}.
\end{tutsteps}
\textbf{What to change next.} Make cancellations proportional to the queue's size (rate $\theta$ per order) and see
how the race changes; make the arrival rates depend on the imbalance, as in the queue-reactive models of One
Quant Book~10.
\end{tutorial}

\section{Build: queues for order books}\label{bld:qm:markov-chains-and-queues:queues}

\begin{build}
\bfield{Purpose} Price the value of a queue position, forecast the next price move from the best queues, and
size message queues in the firm's systems.
\bfield{Interface} \texttt{generator(rates)}; \texttt{stationary(Q)}; \texttt{jump\_chain(Q)};
\texttt{bd\_hitting\_probability(birth, death, n)}; \texttt{bd\_expected\_exit\_time(birth, death, n)};
\texttt{depletion\_cdf(birth, death, start, t)}; \texttt{fill\_time\_cdf(k, nu, mu, t)}; \texttt{race(cdf\_a,
cdf\_b, t)}; \texttt{thomas(a, b, c, d)}.
\bfield{Rules} Rates per second, queues in orders; birth and death rates may be functions of the state for the
first-step solvers; uniformisation, not matrix exponentials, for transient laws; every solver cross-checked by
simulation in the tests.
\bfield{Acceptance tests} \texttt{code/firm/queues/tests/}: gambler's ruin and $i(n - i)/(\lambda + \mu)$ exit times in the
symmetric case; M/M/1 stationary law; depletion law against simulation; race symmetric at equal queues.
\bfield{Stretch} State-dependent rates in the race (a two-dimensional chain solved by Gauss--Seidel); queues fed
by the Hawkes flow of chapter~7.
\end{build}

\begin{omsources}
\item R.~Cont, S.~Stoikov and R.~Talreja, ``A stochastic model for order book dynamics'', \emph{Operations
Research} 58, 2010.
\item R.~Cont and A.~de Larrard, ``Price dynamics in a Markovian limit order market'', \emph{SIAM Journal on
Financial Mathematics} 4, 2013.
\item J.~D.~C.~Little, ``A proof for the queuing formula $L = \lambda W$'', \emph{Operations Research} 9, 1961.
\item D.~G.~Kendall, ``Stochastic processes occurring in the theory of queues and their analysis by the method
of the imbedded Markov chain'', \emph{Annals of Mathematical Statistics} 24, 1953.
\end{omsources}

\section{Exercises}

\begin{exercise}[$\star$]\label{exo:qm:markov-chains-and-queues:1}
In the spread chain of \cref{ex:qm:markov-chains-and-queues:spread}, what is the mean time spent at each spread
before a change, and where does the chain go when it leaves 3 ticks?
\end{exercise}

\begin{exercise}[$\star$]\label{exo:qm:markov-chains-and-queues:2}
An \omterm{def:qm:markov-chains-and-queues:mm1}{M/M/1 queue} serves one message a microsecond and receives 0.8 a microsecond. What are the mean number of
messages in the system and the mean time a message spends there?
\end{exercise}

\begin{exercise}[$\star$]\label{exo:qm:markov-chains-and-queues:3}
A queue of 20 orders gains one at rate 0.9 and loses one at rate 1. What is the probability it reaches 40 before
emptying?
\end{exercise}

\begin{exercise}[$\star\star$]\label{exo:qm:markov-chains-and-queues:4}
Is the spread chain reversible? Check \omterm{def:qm:markov-chains-and-queues:balance}{detailed balance} between 1 and 2 ticks, and Kolmogorov's criterion on
the cycle $1 \to 2 \to 3 \to 1$.
\end{exercise}

\begin{exercise}[$\star\star$]\label{exo:qm:markov-chains-and-queues:5}
For the queue of \cref{exo:qm:markov-chains-and-queues:3} with no upper limit, what are the mean and the median
time to empty, and the probability it empties within a minute? Why do mean and median differ so much?
\end{exercise}

\begin{exercise}[$\star\star$]\label{exo:qm:markov-chains-and-queues:6}
A price level holds 400 lots on average, and 50 lots a second leave it by execution or cancellation. How long does
a lot rest there on average?
\end{exercise}

\begin{exercise}[$\star\star\star$]\label{exo:qm:markov-chains-and-queues:7}
\emph{Coding.} Check the fill probability of the weekend problem by simulating the race event by event, and
compare with the uniformisation result.
\end{exercise}

\begin{exercise}[$\star\star\star$]\label{exo:qm:markov-chains-and-queues:8}
\emph{Find the flaw.} ``Our passive orders get filled 95\% of the time when the book is this imbalanced, so we
join the back of the heavy side whenever we see it.''
\end{exercise}

\section{Problem: The Back of the Queue}

\begin{problem}[{Weekend problem --- 800 lots ahead, 200 on the other side}]\label{pb:qm:markov-chains-and-queues:1}
A market maker's order joins the back of a best bid of 800 lots; the best ask holds 200. Orders are 10 lots.
On each side limit orders join at 0.9 per second, market orders arrive at 0.6 per second and cancellations at
0.4, so orders leave at $\nu = 1$ per second. Take the queues independent.

\medskip\noindent\textbf{Part I --- The model.}
\begin{enumerate}
\item How many orders are ahead of ours, and on the ask?
\item What is the expected time until our order is filled, ignoring the ask?
\item What are the mean and the median time for the ask queue to empty?
\item What is the probability that the ask empties within a minute?
\item What is the probability that it empties at all?
\end{enumerate}

\medskip\noindent\textbf{Part II --- The race.}
\begin{enumerate}[resume]
\item What is the probability that our order fills before the ask empties?
\item What does a simulation of 20\,000 races give?
\item And if the ask held 400 or 800 lots?
\item And if only 20 orders were ahead of ours?
\item What is the probability that the next move of the mid is up, and with equal queues?
\end{enumerate}

\medskip\noindent\textbf{Part III --- First steps.}
\begin{enumerate}[resume]
\item What is the probability that the ask queue reaches 40 orders before it empties?
\item What is the expected time until it first reaches 0 or 40?
\item How would cancellations proportional to the queue size change the drift?
\item What would clustered order flow (chapter~7) do to the tail of the emptying time?
\item What does independence of the two queues leave out?
\end{enumerate}

\medskip\noindent\textbf{Part IV --- Judgement.}
\begin{enumerate}[resume]
\item What is a queue position worth, in this model?
\item When the order does fill in this configuration, what does the fill tell the market maker?
\item Should it join the bid at all?
\item State the \emph{named result}: the fill probability and the expected wait.
\item In one sentence: what decides whether a passive order at the back of a queue fills?
\end{enumerate}
\end{problem}

\section{Interview questions}

\begin{interviewq}[$\star$ \iqroles{researcher, developer}]\label{iq:qm:markov-chains-and-queues:1}
A server handles 10\,000 requests a second and receives 9\,000. On average, how many are waiting and for how long?
\end{interviewq}

\begin{interviewq}[$\star$ \iqroles{researcher}]\label{iq:qm:markov-chains-and-queues:2}
What is the generator of a \omterm{def:qm:markov-chains-and-queues:ctmc}{continuous-time Markov chain}, and how do you get its \omterm{def:qm:stochastic-differential-equations:kolmogorov}{stationary distribution}?
\end{interviewq}

\begin{interviewq}[$\star\star$ \iqroles{trader, researcher}]\label{iq:qm:markov-chains-and-queues:3}
The bid has 800 lots and the ask 200. Which way is the next tick more likely, and why?
\end{interviewq}

\begin{interviewq}[$\star\star$ \iqroles{researcher}]\label{iq:qm:markov-chains-and-queues:4}
A fair random walk starts at 3. What are the probability of hitting 10 before 0 and the expected number of steps?
\end{interviewq}

\begin{interviewq}[$\star\star$ \iqroles{trader}]\label{iq:qm:markov-chains-and-queues:5}
Why do market makers value a place at the front of the queue, and how would you put a number on it?
\end{interviewq}

\begin{interviewq}[$\star\star\star$ \iqroles{developer, researcher}]\label{iq:qm:markov-chains-and-queues:6}
How do you compute the probability that a chain is in a given state after time $t$, for a chain with thousands
of states, stably?
\end{interviewq}
